\chapter{Final Audit: What This Book Actually Established}
\label{ch:final-audit}

This chapter derives nothing. Chapter~\ref{ch:assumption-audit} took a
snapshot after Part II; this chapter classifies the \emph{entire} book,
statement by statement, into four categories:

\begin{description}
\item[{[D] Derived}] — proved within this book's reconstruction, from
admissibility constraints already established, with no new physical
assumption.
\item[{[R] Equivalent reformulation}] — restates an earlier result in
different language; adds no independent content.
\item[{[A] Independently assumed}] — a genuine physical or operational
principle introduced because the admissibility logic alone did not force
it, flagged explicitly at the point of introduction.
\item[{[U] Unresolved}] — an open question this book raises but does not
answer.
\end{description}

The purpose is falsifiability about the book's own achievement: a reader
should be able to check every claim's category against where it was proved,
assumed, or left open, rather than take "quantum mechanics has been
reconstructed" on faith.

\section{Single-System Foundations (Part I)}

\begin{itemize}
\item[{[D]}] Probability simplex, stochastic maps, deterministic maps as
extreme points (Ch.~\ref{ch:simplex}).
\item[{[D]}] Birkhoff--von Neumann theorem; Hardy--Littlewood--P\'olya
majorization; entropy non-decrease as a structural corollary
(Ch.~\ref{ch:birkhoff}).
\item[{[D]}] $n=2$ orthostochastic/unistochastic degeneracy; $n=3$
incompleteness via the Hadamard-matrix obstruction (Ch.~\ref{ch:unistochastic}).
\item[{[D]}] Probability-level composition is not a homomorphism of $\pi$
(Ch.~\ref{ch:composition}).
\item[{[U]}] Whether interference alone selects $\C$ over $\R$ — explicitly
answered \emph{no} within Ch.~\ref{ch:composition} itself; resolved (as a
negative result) rather than left open.
\item[{[D]}] Exact frame closure around a cycle; the general graph
realizability theorem (spanning tree + fundamental-cycle checks)
(Ch.~\ref{ch:context-graphs}).
\item[{[D]}] Triangle-gluing theorem and the real-MUB corollary
(Ch.~\ref{ch:real-mub-obstruction}) — the single-system obstruction to $\R^2$,
independent of composite systems.
\item[{[U]}] Whether every non-orthostochastic unistochastic matrix fails to
be orthostochastic for the \emph{same} structural reason as the $n=3$ DFT
example, or whether distinct mechanisms are at play in general
(Ch.~\ref{ch:unistochastic}).
\end{itemize}

\section{System Composition (Part II)}

\begin{itemize}
\item[{[R]}] Minimal-witness reformulation, $m_\beta(\R,\C;2)=1$ at $K_3$
(Ch.~\ref{ch:minimal-witnesses}) — restates Ch.~\ref{ch:real-mub-obstruction}
in extremal language.
\item[{[U]}] Higher-$d$ minimal witnesses; any quaternionic minimal witness;
forbidden-subgraph characterizations; density thresholds
(Ch.~\ref{ch:minimal-witnesses}, \S\ref{sec:open-extremal}) — an explicit
research programme, not results.
\item[{[A]}] Local tomography, as a principle imposed and tested rather than
derived from Part I (Ch.~\ref{ch:local-tomography}).
\item[{[D]}] The dimension identity distinguishing $\R,\C,\Hset$ under local
tomography, given that principle (Ch.~\ref{ch:local-tomography}).
\item[{[D]}] Tensor products of local bases span the composite space for
$\C$; the concrete missing $\R$ direction $J\otimes J$
(Ch.~\ref{ch:tensor-product}).
\item[{[D]}] Positivity derived operationally from nonnegative test
probabilities; the $|c|\le1$ bound (Ch.~\ref{ch:positivity}).
\item[{[D]}] Partial trace defined and shown to reproduce local statistics;
$J\otimes J$ invisible to both marginals (Ch.~\ref{ch:partial-trace}).
\item[{[D]}] Field-relative separability: the $J\otimes J$ witness for $\R$;
the explicit $\C$-separable decomposition for the entire range
(Ch.~\ref{ch:entanglement}).
\item[{[U]}] Multipartite entanglement, entanglement measures, $d>2$
composite systems in general; what "separable" even means for $\Hset$, given
Ch.~\ref{ch:tensor-product}'s overcompleteness (Ch.~\ref{ch:entanglement}).
\end{itemize}

\section{Recovered Formalism (Part III)}

\begin{itemize}
\item[{[D]}] General-$d$ context-transition identity (Prop.
\ref{prop:general-context-transition}) — a new, easy generalization of the
real 2D case, not previously written down explicitly.
\item[{[A]}] Pure states as frame-completable vectors (Assumption
\ref{ass:frame-completion}, Ch.~\ref{ch:born-rule}) — a genuine new
operational assumption identifying "a state" with a distinguished frame
coordinate.
\item[{[D]}] Pure-state Born rule, given that assumption; density-operator
Born rule, at no further cost beyond Part II's mixture structure
(Ch.~\ref{ch:born-rule}).
\item[{[D]}] The Hermitian/real-sharp-measurement equivalence (one direction
by construction, the other by the cited spectral theorem); the eigenvalue
rule as a corollary (Ch.~\ref{ch:observables}).
\item[{[A]}] Real-valued outcome labels (Remark~\ref{rem:real-labels-assumption})
— flagged explicitly as a modeling choice, not derived.
\item[{[U]}] The spectral theorem itself — cited from classical linear
algebra, not proved within this book's reconstruction
(Ch.~\ref{ch:observables}). Compatible observables and commutators.
\item[{[A]}] Continuous reversible time evolution as a strongly continuous
unitary group (Assumption~\ref{ass:continuous-evolution},
Ch.~\ref{ch:dynamics}) — imported, not derived from Parts I--II's discrete
framework.
\item[{[D]}] The finite-dimensional generator theorem and Schr\"odinger's
equation as its derivative, given that assumption (Ch.~\ref{ch:dynamics}).
\item[{[U]}] That the generator $H$ is energy, or that $\hbar$ is Planck's
constant — flagged explicitly as requiring additional physical
identification not attempted (Ch.~\ref{ch:dynamics}). Open-system/non-unitary
dynamics and the Kraus formalism.
\item[{[D]}] Naimark dilation for finite POVMs, proved constructively
(Ch.~\ref{ch:povm}).
\item[{[U]}] Continuum-outcome POVMs; uniqueness or physical distinguishedness
of a given dilation; instruments and post-measurement state updates
(Ch.~\ref{ch:povm}).
\end{itemize}

\section{A Note on Recent Literature: Representation Versus Structure}
\label{sec:recent-literature}

Two papers appearing during the writing of this book bear directly on the
Kronecker-product convention used throughout Chapters~\ref{ch:local-tomography}
and~\ref{ch:tensor-product}, and are worth citing precisely rather than
leaving this book silent about.

Hoffreumon and Woods~\cite{hoffreumon-woods-2025} construct a real-number
quantum theory reproducing every prediction of standard complex quantum
theory, responding to the 2021 falsification argument of Renou et
al.~\cite{renou2021}. Their construction replaces the Kronecker product with
an alternative representation $\otimes_r$ of the \emph{abstract} tensor
product, satisfying $\Gamma(H \otimes L) = \Gamma(H) \otimes_r \Gamma(L)$ for
a realification map $\Gamma$ decomposing a Hermitian matrix $H = H^{\mathrm{
Sy}} + iH^{\mathrm{An}}$ into real symmetric and antisymmetric parts. The
real matrix standing in for $i$ in their construction is
\begin{equation}
J = \begin{pmatrix} 0 & -1 \\ 1 & 0 \end{pmatrix},
\end{equation}
identical to the $J$ of Chapters~\ref{ch:tensor-product}--\ref{ch:entanglement}.
Barrios Hita, Trushechkin, Kampermann, Epping, and Bru\ss~\cite{barrios-hita-2026}
reach the same empirical conclusion via a different construction, quotienting
the ordinary real tensor product by the phase-redistribution redundancy
$i|\psi_A\rangle \otimes (-i)|\psi_B\rangle = |\psi_A\rangle\otimes
|\psi_B\rangle$ carries in the complex theory.

\begin{remark}[What changes, and what does not]
Both papers confirm that the Kronecker product is one admissible matrix
representation of composition, selected by convention rather than forced by
locality or linearity alone — precisely the distinction this book's
admissibility methodology asks about everywhere else, and had not asked about
here. Chapters~\ref{ch:local-tomography} and~\ref{ch:tensor-product} used the
Kronecker product throughout without flagging it as a choice; every claim
proved there remains correct \emph{for that representation}, but the audit
should not be read as having shown Kronecker composition to be the unique or
necessary choice.

This point is older than either paper, and worth stating precisely rather
than leaving implicit. That a complex number is equivalent to a real
$2\times2$ rotation-scaling matrix, $a+bi \leftrightarrow \begin{psmallmatrix}
a&-b\\b&a\end{psmallmatrix}$ — Needham's \emph{amplitwist}: multiplication
by $z=re^{i\theta}$ as simultaneous scaling by $r$ and rotation by
$\theta$~\cite{needham1997} — has nothing to do with quantum mechanics and
predates these papers by decades. What is genuinely new in
Hoffreumon--Woods and Barrios Hita et al.\ is not this realification (classical),
but the demonstration that a full \emph{multipartite} composition law can be
built on top of it while preserving locality and reproducing every quantum
prediction — a nontrivial claim about composition, not merely a restatement
of the classical fact that $i$ is a rotation-scaling operator in disguise.
The two should not be conflated: the elementary observation licenses no
conclusion about quantum composition on its own, and the papers' technical
content is precisely the part the elementary observation does not supply.

What does \emph{not} change is the assumption count of \S\ref{sec:assumptions-collected}.
Local tomography is not derived by either paper from weaker principles;
Hoffreumon and Woods obtain it for their real theory via a one-to-one
correspondence with complex quantum theory's own local tomography, which is
a transfer argument, not a derivation from representation locality alone. The
distinct question these papers answer — whether complex \emph{scalars} are
necessary, given that some coherent structure equivalent to $J^2=-I$ survives
into the composition law regardless — is a different, and in a precise sense
narrower, claim than whether the admissibility constraints of Parts I--II
force $\C$. This book's real-MUB and local-tomography arguments (Chapters
\ref{ch:real-mub-obstruction}, \ref{ch:local-tomography}) are unaffected by
either paper: they rule out $\R$ under the Kronecker representation
specifically, and whether they still rule out $\R$ under a
$\otimes_r$-style representation is a genuinely open question this book does
not address.
\end{remark}

\begin{ledgerentry}[Addendum]
\textbf{Not established:} whether Chapters~\ref{ch:local-tomography} and
\ref{ch:tensor-product}'s results survive replacement of the Kronecker
product with a non-Kronecker representation of composition in the style of
Hoffreumon--Woods or Barrios Hita et al. This is left as an explicit open
question rather than assumed away, added to the record of
\S\ref{sec:open-extremal} and \S\ref{sec:what-remains}.
\end{ledgerentry}

\section{The Assumptions, Collected}
\label{sec:assumptions-collected}

Stripped of derivation and restated together, this is the complete list of
what this book's reconstruction has \emph{added} to bare admissibility
reasoning, beyond what Parts I--II derive from constraint alone:

\begin{enumerate}
\item Local tomography (Ch.~\ref{ch:local-tomography}) — tested, not forced.
\item Pure states are frame-completable vectors (Ch.~\ref{ch:born-rule}).
\item Measurement outcomes are labeled by real numbers (Ch.~\ref{ch:observables}).
\item Time evolution is a strongly continuous unitary group
(Ch.~\ref{ch:dynamics}).
\end{enumerate}

Four assumptions. Everything else in the recovered formalism — the Born rule
in all its forms, the Hermitian-observable correspondence, the eigenvalue
rule, Schr\"odinger's equation, and the Naimark dilation of POVMs — is
derived from these four plus Parts I--II's admissibility structure, not
independently postulated. Whether four is a small number is a matter of
comparison the reader can make directly against
Sakurai--Napolitano's~\cite{sakurainapolitano} or Gross's~\cite{gross}
postulate lists, cited throughout this book precisely so that
comparison is possible.

\begin{scopebox}
\scopetitle{}A small number of assumptions is not the same claim as zero
assumptions, and it is not the same claim as inevitability. Four
independently chosen principles, none derived from anything more
primitive, is compatible with several different theories of quantum
mechanics being equally well-founded — it is not a proof that this
particular four is the unique starting point physics could have had.
Nothing in this count should be read as "quantum mechanics is the only
way the pieces fit together"; it should be read as "this is exactly what
standard quantum mechanics costs, on this specific route, stated
precisely enough to compare against other routes" — including the
routes discussed in \S\ref{sec:recent-literature} below, which reach
empirically equivalent theories via a different, likewise small,
assumption set.
\end{scopebox}

\begin{ledgerentry}[Closing entry]
This chapter adds no new mathematics, with one exception: \S\ref{sec:recent-literature}
notes that the Kronecker-product convention used in Chapters
\ref{ch:local-tomography}--\ref{ch:tensor-product} is a representation
choice, not a forced one, per recent literature — itself an audit-relevant
fact, not a new derivation. Otherwise: every derived statement traces to a
cited line, every assumption is numbered above, and every unresolved
question is listed rather than quietly dropped. Any future chapter that
treats an [A]-tagged or [U]-tagged item as though it were [D] should be read
as an error against this audit, not a new result.
\end{ledgerentry}
